\chapter{Projekt i wykorzystane technologie}
\section{Wprowadzenie}
W celu dopasowania oprogramowania do założeń przedstawionych w poprzednim rozdziale, do realizacji wybrano narzędzia pozwalające na bezbolesną integrację, możliwie najniższe wymagania sprzętowe, niską ilość zależności (ang. dependencies) oraz przejrzystość kodu źródłowego. W tym rozdziale przedstawiona zostanie ogólna architektura rozwiązania, jego elementy, oraz technologie wykorzystane do jego stworzenia.

\section{Architektura rozwiązania}
% Diagram ogólny
% Czy nie powinniśmy opisać konfigu? może gdzieś później opisać w jaki sposób program startuje? 
% co wgl z opisem deploymentu? jak możemy zdeployować tą appkę na czyiś komputer, albo gdziekolwiek? docker?
\subsection{Wejście}
Proces rozpoczyna się poprzez przyjęcie danych w formacie JSON % citation
od systemu zewnętrznego poprzez żądanie POST % citation
do endpointu \texttt{/invoices}, eksponowanego przez aplikację. 
Format JSON został wybrany ze względu na swoją kompatybilność z HTTP API oraz szerokie wsparcie przez systemy zintegrowane i aplikacje backendowe. Jest on również czytelny dla człowieka. Aplikacja akceptuje dane w formacie JSON, służącym jako pośredni, integracyjny format, i używa ich do stworzenia bardziej ustrukturyzowanego pliku XML dostosowanego do oficjalnego schematu (FA3) udostępnionego przez Ministerstwo Finansów
% citation
.

\noindent Oprogramowanie spodziewa się następujących danych:
\begin{small}
\begin{longtable}{>{\raggedright\arraybackslash}p{4cm} >{\raggedright\arraybackslash}p{2cm} >{\raggedright\arraybackslash}p{7cm}}
\caption{Pola żądania wejściowego HTTP przekazywanego do endpointu \texttt{/invoices}}
\label{tab:invoice_input_fields} \\
\toprule
\textbf{Pole} & \textbf{Wymagalność} & \textbf{Opis} \\
\midrule
\endfirsthead

\toprule
\textbf{Pole} & \textbf{Wymagalność} & \textbf{Opis} \\
\midrule
\endhead

\bottomrule
\endfoot

invoice\_type & TAK & Typ faktury - VAT, korygująca, zaliczkowa, rozliczeniowa, uproszczona, korygująca rozliczeniowa lub korygująca zaliczkowa. \\
invoice\_number & TAK & Numer faktury w systemie zewnętrznym. \\
issue\_date & TAK & Data wystawienia faktury w formacie RRRR-MM-DD. \\
currency & TAK & Kod waluty zgodny z ISO 4217. \\
exchange\_rate & TAK / NIE & Kurs waluty dla walut innych niż PLN. \\
total\_amount & TAK & Łączna wartość brutto faktury. \\
seller & TAK & Dane sprzedawcy - NIP, nazwa oraz adres. \\
buyer & NIE & Dane nabywcy - niewymagane w przypadku faktury uproszczonej. \\
items & TAK & Lista wierszy faktury, zawierająca nazwę, ilość, wartość netto oraz stawkę podatku dla każdej usługi/produktu. \\
tax\_breakdowns & TAK & Wartości netto i podatku, podzielone przez stawki VAT. \\
payment & NIE & Dane dotyczące płatności - termin płatności, metoda i rachunek bankowy. \\
flags & NIE & Dodatkowe znaczniki, przykładowo płatność podzielona lub metoda kasowa. \\
callback\_url & NIE & Adres URL, na który oprogramowanie wyśle informację o statusie przesyłki. \\
original\_invoice\_number & NIE & Numer oryginalnej faktury, wymagany w przypadku faktur korygujących. \\
original\_issue\_date & NIE & Data wystawienia oryginalnej faktury, wymagana w przypadku faktur korygujących. \\
original\_ksef\_id & NIE & Numer identyfikacyjny faktury w KSeF, wymagany w przypadku faktur korygujących. \\
correction\_reason & NIE & Powód korekty, wymagany w przypadku faktur korygujących. \\
corrected\_buyer & TAK / NIE & Poprawione dane nabywcy, wymagane w przypadku faktur korygujących. \\
\end{longtable}
\end{small}

\subsection{Walidacja otrzymanych danych}
Treść żądania jest następnie przetwarzana - występuje walidacja podstawowych danych faktury w celu stwierdzenia, czy powinna zostać podjęta próba przesłania jej do Krajowego Systemu e-Faktur. Sprawdzane dane to:
\begin{itemize}
    \item Typ faktury - faktura VAT, faktura korygująca, faktura zaliczkowa, faktura rozliczeniowa, faktura uproszczona, faktura korygująca rozliczeniowa lub faktura korygująca zaliczkowa.
    \item W przypadku faktur korygujących, sprawdzane są pola dotyczące danych korekty - powód, data wystawienia oryginalnej faktury, poprawność daty oraz poprawione dane kontrahenta. W przypadku faktur niekorygujących, pola sprawdzane są z założeniem, że powinny być puste.
    \item Podstawowa walidacja waluty - ilość znaków, oraz w przypadku walut innych niż PLN, kurs między określoną walutą a PLN. 
    \item Wiersze sprzedaży - aplikacja sprawdza, czy faktura nie jest pusta. Następnie waliduje każdy wiersz sprzedaży, poprzez sprwadzenie ilości określonego produktu/usługi, ceny na jednostkę oraz nałożenia odpowiedniego podatku. W przypadku faktur niekorygujących, odrzuca wartości netto poniżej zera. Wartości netto oraz podatku są zsumowane i walidowane przeciwko całkowitej wartości faktury.
    \item Dane kontrahentów - sprawdzane są Numery Identyfikacji Podatkowej obu stron wykonanej transakcji poprzez sprawdzenie ich odpowiednim algorytmem % przypis!!!
        oraz nazwy i adresy stron. Walidowany jest również kod kraju sprzedawcy, a w przypadku jego braku, następuje wprowadzenie wartości domyślnej - PL. W przypadku faktury uproszczonej, dane kupującego nie są wymagane, zgodnie z ustawą % PODAJ DANE USTAWY .
\end{itemize}

\subsection{Utworzenie pliku w formacie XML}
% my nie mamy nawet wsparcia dla grup vat..... strasznie duzo magicznych liczb w transformacji do xml...
Dane są następnie transformowane do formatu XML, poprzez mapowanie odpowiednich pól JSON do pól wymaganych przez oficjalny schemat XSD. 

\subsection{Walidacja danych w formacie XML}
Aplikacja waliduje nowo utworzone dane XML wobec schematu XSD
% przypis co to XSD?
udostępnionego przez Ministerstwo Finansów.

\subsection{Zapis danych}
W przypadku pomyślności poprzednich kroków, dane zapisywane są w lokalnej bazie danych SQLite.

\noindent Zapisywane są następujące dane:
\begin{small}
\begin{longtable}{>{\raggedright\arraybackslash}p{4cm} >{\raggedright\arraybackslash}p{2cm} >{\raggedright\arraybackslash}p{7cm}}
\caption{Najważniejsze kolumny tabeli Invoices w lokalnej bazie danych}
\label{tab:sqlite_invoices} \\
\toprule
\textbf{Kolumna} & \textbf{Typ} & \textbf{Opis} \\
\midrule
\endfirsthead

\toprule
\textbf{Kolumna} & \textbf{Typ} & \textbf{Opis} \\
\midrule
\endhead

\bottomrule
\endfoot

id & INTEGER & Klucz główny rekordu. \\
external\_id & TEXT & Numer faktury w systemie zewnętrznym. \\
raw\_json & TEXT & Oryginalne dane otrzymane w formacie JSON. \\
raw\_xml & TEXT & Faktura przekonwertowana do formatu XML. \\
status & TEXT & Status przetwarzania faktury - PENDING, SENT, FAILED lub UNKNOWN. \\
ksef\_id & TEXT & Numer faktury w KSeF. \\
ksef\_error & TEXT & Błąd zwrócony przez KSeF w przypadku odrzucenia. \\
upo\_xml & TEXT & Urzędowe Poświadczenie Odbioru w formacie XML. \\
attempt\_count & INTEGER & Liczba prób wysłania faktury do KSeF. \\
created\_at & DATETIME & Data utworzenia. \\
updated\_at & DATETIME & Data ostatniej aktualizacji. \\
submission\_reference & TEXT & Identyfikator przesłania, nadany przez KSeF. \\
callback\_url & TEXT & Opcjonalny adres do przesłania informacji dotyczącej statusu faktury, po zakończeniu przerwarzania. \\
webhook\_delivered & BINARY & Informacja, czy dane dotyczące statusu przesyłki zostały zwrócone. \\
webhook\_attempt\_count & INTEGER & Liczba prób dostarczenia informacji zwrotnej. \\
webhook\_error & TEXT & Błąd w przypadku niepomyślnego dostarczenia informacji zwrotnej. \\
\end{longtable}
\end{small}

\subsection{Lokalna kolejka oraz wysyłka}
Aplikacja skanuje bazę danych w poszukiwaniu faktur oznaczonych jako niewysłane oraz dodaje je do lokalnej kolejki. Jeśli kolejka nie jest pusta, otwierana jest sesja interaktywna przez autentykację w KSeF, faktury są wysyłane oraz monitorowane pod względem statusu wysyłki. W przypadku akceptacji, faktura oznaczana jest jako wysłana - w przypadku nieznanego statusu, aplikacja ponowi zapytanie dotyczące faktury. Jeśli zaś dojdzie do porażki, aplikacja spróbuje wysłać fakturę ponownie w następnym cyklu działania, do momentu gdy nie zostanie osiągnięty określony przez użytkownika w pliku konfiguracyjnym limit prób. Wyjątkiem jest tutaj błąd połączenia z internetem - w tym wypadku, faktura pozostanie w kolejce, do momentu, gdy możliwa będzie wysyłka.

\subsection{Webhooki}
Jeśli w oryginalnych danych przekazanych aplikacji żądaniem figurował adres, na który wysłana powinna zostać informacja zwrotna dotycząca statusu wysyłki, mechanizm webhooków (z ang. haków sieciowych - aplikacja "zahacza" o dostawcę danych, na potrzebę przyszłej interakcji)
prześle zewnętrznemu systemowi dane w formacie JSON po oznaczeniu faktury jako wysłana lub odrzucona, potwierdzając pomyślność operacji lub podając powód odrzucenia faktury. W przypadku błędu w trakcie próby przekazania danych, aplikacja podejmie kolejne próby, aż do osiągnięcia określonego przez użytkownika limitu prób.

\noindent Informacja zwrotna przesyłana jest w następującej postaci:
\begin{small}
\begin{longtable}{>{\raggedright\arraybackslash}p{4cm} >{\raggedright\arraybackslash}p{2cm} >{\raggedright\arraybackslash}p{7cm}}
\caption{Pola żądania HTTP z informacją zwrotną przesyłanego do systemu zewnętrznego}
\label{tab:webhook_payload} \\
\toprule
\textbf{Pole} & \textbf{Typ} & \textbf{Opis} \\
\midrule
\endfirsthead

\toprule
\textbf{Pole} & \textbf{Typ} & \textbf{Opis} \\
\midrule
\endhead

\bottomrule
\endfoot

event & TEXT & Typ zdarzenia, przykładowo invoice.accepted lub invoice.failed. \\
invoice\_id & INTEGER & Identyfikator faktury w aplikacji - klucz główny w bazie danych. \\
external\_id & TEXT & Numer faktury pochodzący z systemu zewnętrznego. \\
status & TEXT & Status wysyłki - SENT lub FAILED. \\
ksef\_reference\_number & TEXT / NULL & Numer faktury w KSeF, jeśli faktura została przyjęta. \\
submission\_reference & TEXT / NULL & Identyfikator przesłania, nadany przez KSeF. \\
error\_message & TEXT / NULL & Błąd, który nastąpił w przypadku niepowodzenia. \\
timestamp & DATETIME & Czas przesyłki informacji zwrotnej. \\
\end{longtable}
\end{small}

\subsection{Panel użytkownika}
Aplikacja eksponuje również prosty panel na endpoincie \texttt{/ui/invoices} umożliwiający wgląd do historii wysyłek, filtrowanie faktur według statusu, wyszukiwanie pojedynczych rekordów oraz ich manualne usuwanie w przypadku niepomyślnej wysyłki i osiągnięcia limitu prób. Panel automatycznie odświeża swoją zawartość. Panel wykonany został z użyciem HTMX % przypis czy opisać HTMX w sekcji technologie?
    oraz ostylowany za pomocą Tailwind CSS.

\section{Komunikacja z KSeF}
Podstawowym elementem działania aplikacji jest interakcja z samym Krajowym Systemem e-Faktur - uwierzytelnienie, otwieranie i zamykanie sesji interaktywnych w celu wysyłki faktur oraz sprawdzanie statusu i odbiór Urzędowych Poświadczeń Odbioru. 

\subsection{Uwierzytelnienie}
Uwierzytelnienie następuje natychmiast po uruchomieniu aplikacji i podtrzymywane jest przez cały okres jej działania. Aplikacja wykonuje puste żądanie POST do endpointu \texttt{/auth/challenge} eksponowanego przez KSeF, w odpowiedzi uzyskując \enquote{wyzwanie} - timestamp, który po połączeniu z pozyskanym wcześniej przez użytkownika tokenem API z panelu KSeF jest szyfrowany z użyciem klucza publicznego automatycznie pobieranego przez aplikację. Przesyłane następnie jest żądanie do endpointu \texttt{/auth/ksef-token}, zwracające tymczasowy token dostępu. Następnie, token ten wysyłany jest do \texttt{/auth/token/redeem}, pozyskując token dostępu oraz drugi, pozwalający na odświeżenie go. Oprogramowanie automatycznie sprawdza, kiedy dojdzie do przedawnienia, i automatycznie odświeży status uwierzytelnienia, jeśli jest taka potrzeba.

\subsection{Otwieranie i zamykanie sesji}
Do otwierania oraz zamykania tzw. sesji interaktywnych, pozwalających na wysyłkę faktur, dochodzi gdy aplikacja rejestruje pojawienie się nowej faktury lub faktur do wysyłki w bazie danych. Sesja otwierana jest poprzez przesłanie żądania z zaszyfrowanym kluczem symetrycznym oraz wektorem inicjalizacji do endpointu \texttt{/sessions/online}, otrzymując w odpowiedzi identyfikator sesji oraz czas jej ważności. Po wysyłce sesja jest zamykana przez przesłanie żądania do endpointu \texttt{/sessions/online/x/close}, gdzie \enquote{x} jest zapisanym przy otwarciu identyfikatorem.

\subsection{Przesyłanie faktur}
Faktury z bazy danych, w formacie XML, są następnie szyfrowane oraz przesyłane do endpointu \texttt{/sessions/online/x/invoices}, gdzie \enquote{x} to identyfikator sesji interaktywnej. Zwracany jest identyfikator przesyłki, który wykorzystywany jest do sprawdzenia statusu wysyłki. Niepotwierdzane jest jednak czy faktura została przyjęta przez KSeF.

\subsection{Sprawdzenie statusu i pobranie Urzędowego Poświadczenia Odbioru}
Aplikacja sprawdza status faktury aż do otrzymania odpowiedzi, lub osiągnięcia limitu ponownych żądań ustalonego przez użytkownika w konfiguracji. W przypadku przyjęcia faktury przez KSeF, aplikacja zapisze Urzędowe Poświadczenie Odbioru w bazie danych, a w przypadku odrzucenia oznaczy w bazie danych porażkę. W przypadku otrzymywania odpowiedzi, że faktura dalej jest przetwarzana, oznaczy jej status jako nieznany, oraz podejmie próbę sprawdzenia jej statusu przy kolejnym cyklu kolejki.

\section{Cykl życia aplikacji}
W tej sekcji opisane będzie ogólne działanie - od konfiguracji i uruchomienia, przez działanie, po mechanizm bezpiecznego zamknięcia oraz w jaki sposób aplikacja jest wdrażana kontenerowo.

\subsection{Konfiguracja}
Oprogramowanie jest konfigurowane za pomocą pliku \enquote{config.yaml}, którego format został wybrany ze względu na swoją prostotę, czytelność dla człowieka oraz powielenie użycia we wdrożeniu kontenerowym (plik \enquote{docker-compose.yaml}), któremu dedykowana jest osobna subsekcja. 


\noindent Struktura pliku konfiguracyjnego jest następująca:

\begin{small}
\begin{longtable}{>{\raggedright\arraybackslash}p{4.8cm} >{\raggedright\arraybackslash}p{1.7cm} >{\raggedright\arraybackslash}p{6.5cm}}
\caption{Plik konfiguracyjny aplikacji}
\label{tab:config_fields} \\
\toprule
\textbf{Parametr} & \textbf{Typ} & \textbf{Opis} \\
\midrule
\endfirsthead

\toprule
\textbf{Parametr} & \textbf{Typ} & \textbf{Opis} \\
\midrule
\endhead

\bottomrule
\endfoot

ksef.nip & TEXT & Numer Identyfikacji Podatkowej podmiotu, który reprezentuje użytkownik. \\
ksef.url & TEXT & Oficjalny adres API KSeF, na wypadek zmiany, lub chęci użytkownika do przetestowania oprogramowania na środowisku testowym. \\
ksef.token & TEXT & Oficjalny token API pozyskany przez użytkownika z Panelu Podatnika KSeF. \\
ksef.http\_timeout\_sec & INTEGER & Maksymalny czas oczekiwania na odpowiedź HTTP z KSeF. \\
ksef.auth\_retry\_delay\_sec & INTEGER & Przerwa w sekundach pomiędzy kolejnymi próbami autentykacji w KSeF. \\
ksef.polling\_delay\_sec & INTEGER & Przerwa w sekundach pomiędzy kolejnymi zapytaniami o status wysłanej do KSeF faktury. \\
ksef.confirmation\_max\_attempts & INTEGER & Maksymalna liczba prób sprawdzenia statusu przesłanej faktury. \\
sqlite.db\_path & TEXT & Ścieżka do lokalnej bazy danych, wybierana przez użytkownika. \\
sqlite.busy\_timeout\_ms & INTEGER & Czas oczekiwania SQLite na zwolnienie blokady bazy danych. \\
server.port & TEXT & Port HTTP, na którym uruchamiana jest aplikacja. \\
user.max\_retries & INTEGER & Maksymalna liczba prób wysyłki faktury lub webhooka w przypadku błędów. \\
xsd\_path & TEXT & Ścieżka do pliku schematu XSD używanego do walidacji faktur w formacie XML. \\
dash\_page\_size & INTEGER & Liczba faktur wyświetlanych na jednej stronie panelu użytkownika. \\
polling\_interval & INTEGER & Czas w sekundach między skanowaniem bazy danych w poszukiwaniu faktur do wysyłki. \\
shutdown\_timeout\_sec & INTEGER & Maksymalny czas oczekiwania na bezpieczne zamknięcie aplikacji. \\
sender\_batch\_size & INTEGER & Maksymalna liczba faktur lub webhooków wysyłanych w jednym cyklu działania aplikacji. \\
sender\_worker\_limit & INTEGER & Limit równoległych wysyłek faktur lub webhooków wykonywanych naraz przez aplikację. \\
\end{longtable}
\end{small}

\subsection{Uruchomienie aplikacji}
Po uruchomieniu aplikacji, tworzone są loggery odpowiadające za rejestrowanie wszystkich działań aplikacji. Otwierany jest wcześniej wspomniany plik konfiguracyjny, a jego wartości są walidowane. Oprogramowanie utworzy odpowiednią bazę danych SQLite, lub spróbuje połączyć się z już isntiejącą, jeśli owa istnieje. Następnie wczytywany jest schemat XSD służący do walidacji faktur w formacie XML oraz tworzony jest specjalny walidator.

Tworzony jest klient HTTP oraz struktura aplikacji, z użyciem parametrów z pliku konfiguracyjnego, oraz uruchamiany jest serwer HTTP, asynchronicznie z pętlą odpowiadającą za skanowanie bazy danych w poszukiwaniu faktur. Przygotowywany jest również kontekst pozwalający później na bezpieczne zamknięcie aplikacji. 

\subsection{Działanie aplikacji}
Aplikacja w ramach swojego działania eksponuje odpowiednie endpointy API HTTP oraz rozpoczyna cykliczne skanowanie bazy danych SQLite.

\noindent Eksponowane są następujące endpointy:

\begin{small}
\begin{longtable}{>{\raggedright\arraybackslash}p{3cm} >{\raggedright\arraybackslash}p{4cm} >{\raggedright\arraybackslash}p{6cm}}
\caption{Endpointy eksponowane przez aplikację}
\label{tab:app_endpoints} \\
\toprule
\textbf{Metoda} & \textbf{Endpoint} & \textbf{Opis} \\
\midrule
\endfirsthead

\toprule
\textbf{Metoda} & \textbf{Endpoint} & \textbf{Opis} \\
\midrule
\endhead

\bottomrule
\endfoot

POST & /invoices & Główny endpoint wejściowy aplikacji. Przyjmuje faktury w formacie JSON. \\
GET & /ui/invoices & Endpoint pozwalający na dostęp do panelu użytkownika. \\
GET & / & Podstawowy endpoint pozwalający na dostęp do panelu użytkownika przez przekierowanie do \texttt{/ui/invoices}. \\
GET & /ui/invoicetable & Endpoint zwracający interfejs wierszy z fakturami, wykorzystywany do odświeżania widoku z użyciem HTMX. \\
GET & /ui/invoice?id=\{id\} & Endpoint zwracający szczegółowy widok pojedynczej faktury. \\
DELETE & /deleteinvoice?id=\{id\} & Endpoint pozwalający na manualne usunięcie faktury. \\
GET & /health/live & Endpoint pozwalający sprawdzić, czy aplikacja działa. \\
GET & /health/ready & Endpoint pozwalający sprawdzić, czy aplikacja jest gotowa do obsługi wysyłanych faktur. \\
\end{longtable}
\end{small}
